\section{Anillos regulares}
\label{section:0.17-es}

\subsection{Definición de los anillos regulares}
\label{subsection:0.17.1-es}

\begin{proposition}[17.1.1]
\label{0.17.1.1-es}
Sean $A$ un anillo local noetheriano de dimensión $n$, $\mathfrak{m}$ su
ideal maximal, $k=A/\mathfrak{m}$ su cuerpo residual. Las condiciones siguientes
son equivalentes~:
\begin{enumerate}
  \item[\rm a)] El homomorfismo canónico sobreyectivo de $k$-módulos graduados
  (15.1.1.1)
  \[
    \varphi:\mathbf{S}^*_k(\mathfrak{m}/\mathfrak{m}^2)
    \longrightarrow\operatorname{gr}^*_{\mathfrak{m}}(A)
  \]
  (donde el segundo miembro es el $k$-módulo graduado asociado a $A$ provisto de
  la filtración $\mathfrak{m}$-preádica) es biyectivo.
  \item[\rm b)] Se tiene
  $\operatorname{rg}_k(\mathfrak{m}/\mathfrak{m}^2)=n$.
  \item[\rm c)] El ideal $\mathfrak{m}$ admite un sistema de generadores de $n$
  elementos.
  \item[\rm d)] El ideal $\mathfrak{m}$ admite un sistema de generadores que es
  una sucesión $A$-regular.
\end{enumerate}
\end{proposition}


\begin{proof}
En virtud del lema de Nakayama,
$\operatorname{rg}_k(\mathfrak{m}/\mathfrak{m}^2)$ es el menor número de
elementos de un sistema de generadores de $\mathfrak{m}$, por lo que b) y c) son
equivalentes. Por otra parte, si $(x_i)_{1\leq i\leq r}$ es un sistema de
generadores de $\mathfrak{m}$ cuyas clases $\overline{x}_i$ módulo
$\mathfrak{m}^2$ forman una base del $k$-espacio vectorial
$\mathfrak{m}/\mathfrak{m}^2$, $\mathbf{S}^*(\mathfrak{m}/\mathfrak{m}^2)$
es isomorfo al anillo de
\oldpage[0\textsubscript{IV-1}]{40}
polinomios $k[\mathrm{T}_1,\ldots,\mathrm{T}_r]$~; habida cuenta de que todo
$A$-módulo de tipo finito es separado para la filtración $\mathfrak{m}$-preádica
($0_{\mathrm I}$, 7.3.5), se deduce de (15.1.9) que las condiciones a) y d)
son equivalentes~; además, como toda sucesión $A$-regular de elementos de
$\mathfrak{m}$ tiene a lo sumo $n$ elementos (16.4.1), se ve que d) implica c).
Queda demostrar que c) implica a).

Pongamos, para abreviar,
$\mathrm{S}=\mathbf{S}^*(\mathfrak{m}/\mathfrak{m}^2)
=k[\mathrm{T}_1,\ldots,\mathrm{T}_n]$,
$\mathrm{G}=\operatorname{gr}^*_{\mathfrak{m}}(A)$, y consideremos la sucesión
exacta
$0\longrightarrow\mathfrak{J}\longrightarrow\mathrm{S}
\xrightarrow{\varphi}\mathrm{G}\longrightarrow0$, donde el núcleo
$\mathfrak{J}$ de $\varphi$ es, por tanto, un ideal graduado de $\mathrm{S}$. Para
todo entero $s>0$, se tiene
\begin{equation}
\label{0.17.1.1.1-es}
\tag{17.1.1.1}
  \binom{s+n-1}{n-1}
  =\operatorname{long}(\mathrm{S}_s)
  =\operatorname{long}(\mathrm{G}_s)
  +\operatorname{long}(\mathfrak{J}_s).
\end{equation}

Supongamos $\mathfrak{J}\neq0$, de modo que existe un elemento homogéneo
$u\in\mathfrak{J}$ de grado $h>0$~; $\mathfrak{J}_s$ contiene
$u\mathrm{S}_{s-h}$ para $s\geq h$, y como $u\neq0$ y $\mathrm{S}$ es
íntegro, $u\mathrm{S}_{s-h}$ es isomorfo a $\mathrm{S}_{s-h}$ como
$k$-espacio vectorial~; se tendría, por tanto,
$\operatorname{long}(\mathfrak{J}_s)\geq
\operatorname{long}(\mathrm{S}_{s-h})=\binom{s-h+n-1}{n-1}$, de donde, para
$s\geq h$,
\begin{equation}
\label{0.17.1.1.2-es}
\tag{17.1.1.2}
  \operatorname{long}(\mathrm{G}_s)
  \leq\binom{s+n-1}{n-1}-\binom{s-h+n-1}{n-1}.
\end{equation}
Ahora bien, el segundo miembro de (17.1.1.2) es un polinomio en $s$ de grado
$\leq n-2$~; pero se tiene
$\mathrm{G}_s=\mathfrak{m}^s/\mathfrak{m}^{s+1}$, y para $s$ suficientemente grande,
$\operatorname{long}(\mathrm{G}_s)$ es un polinomio en $s$ de grado exactamente
igual a $n-1$ cuyo coeficiente dominante es $>0$ (16.2.1), lo que contradice
la desigualdad (17.1.1.2). C.Q.F.D.
\end{proof}

\begin{definition}[17.1.2]
\label{0.17.1.2-es}
Se dice que un anillo local noetheriano que verifica las condiciones equivalentes de
(17.1.1) es \emph{regular}.
\end{definition}

\begin{corollary}[17.1.3]
\label{0.17.1.3-es}
Un anillo local regular es íntegro, íntegramente cerrado y un anillo de
Cohen-Macaulay.
\end{corollary}

\begin{proof}
La última aserción se deduce de (17.1.1, d))~; por otra parte, si $A$ es
regular, $\operatorname{gr}^*_{\mathfrak{m}}(A)$ es íntegro y completamente
íntegramente cerrado, por ser isomorfo a $k[\mathrm{T}_1,\ldots,\mathrm{T}_n]$~;
se concluye que $A$ posee también estas dos propiedades (Bourbaki,
\emph{Alg.~comm.}, cap.~III, \S~2, no~3, cor. de la prop.~1 y cap.~V,
\S~1, no~5, prop.~15).
\end{proof}

\begin{examples}[17.1.4]
\label{0.17.1.4-es}
\begin{enumerate}
  \item[\rm (i)] Un anillo local regular de \emph{dimensión} $0$, por ser
  íntegro en virtud de (17.1.3), es necesariamente un \emph{cuerpo}, y recíprocamente.
  \item[\rm (ii)] Para que un anillo local noetheriano de \emph{dimensión} $1$
  sea regular, es necesario y suficiente que sea un \emph{anillo de valuación
  discreta}~: en efecto, decir que un anillo local noetheriano de dimensión $1$ es
  regular significa que su ideal maximal es \emph{principal}, y la conclusión
  se deduce de Bourbaki, \emph{Alg.~comm.}, cap.~VI, \S~3, no~6, prop.~9.
  \item[\rm (iii)] Sea $k$ un cuerpo~; el anillo de series formales
  $A=k[[\mathrm{T}_1,\ldots,\mathrm{T}_n]]$ es un \emph{anillo regular de
  dimensión $n$}~; en efecto, es claro que los $\mathrm{T}_i$
  $(1\leq i\leq n)$ generan el ideal maximal $\mathfrak{m}$ de $A$ y
  forman una sucesión $A$-regular, pues
  $A/(\mathrm{T}_1A+\cdots+\mathrm{T}_iA)$ es isomorfo a
  $k[[\mathrm{T}_{i+1},\ldots,\mathrm{T}_n]]$.
\end{enumerate}
\end{examples}

\begin{proposition}[17.1.5]
\label{0.17.1.5-es}
Para que un anillo local noetheriano $A$ sea regular, es necesario y suficiente que
su completado $\widehat{A}$ lo sea.
\end{proposition}

\begin{proof}
\oldpage[0\textsubscript{IV-1}]{41}
En efecto, el ideal maximal de $\widehat{A}$ es $\mathfrak{m}\widehat{A}$ y se
sabe que $\mathfrak{m}^h/\mathfrak{m}^{h+1}$ y
$(\mathfrak{m}\widehat{A})^h/(\mathfrak{m}\widehat{A})^{h+1}$ son
$k$-espacios vectoriales isomorfos~; la condición a) de (17.1.1) es, por tanto, la
misma para $A$ y $\widehat{A}$.
\end{proof}

\begin{definition}[17.1.6]
\label{0.17.1.6-es}
Dado un anillo local noetheriano $A$, de ideal maximal $\mathfrak{m}$, se
llama \emph{sistema regular de parámetros de $A$} a un sistema de parámetros
para $A$ (16.3.5) que genera $\mathfrak{m}$.
\end{definition}

La existencia de tal sistema equivale a que $A$ sea regular (17.1.1)~;
si $(x_i)_{1\leq i\leq n}$ es un sistema regular de parámetros, es un
sistema minimal de generadores de $\mathfrak{m}$, cuyas clases módulo
$\mathfrak{m}^2$ forman una base de $\mathfrak{m}/\mathfrak{m}^2$ sobre $k$~;
en virtud de (17.1.1, a)) y de (15.1.9), $(x_i)$ es una sucesión $A$-regular
maximal (de donde procede la terminología).

Se advertirá, sin embargo, que, en un anillo regular, un sistema de
parámetros para $A$ que sea también una sucesión $A$-regular \emph{no es
necesariamente} un sistema regular de parámetros, como muestra el ejemplo
de las potencias $k$-ésimas de un sistema regular de parámetros (15.1.20).

\begin{proposition}[17.1.7]
\label{0.17.1.7-es}
Sean $A$ un anillo local noetheriano, $\mathfrak{m}$ su ideal maximal,
$k=A/\mathfrak{m}$ su cuerpo residual, $(x_i)_{1\leq i\leq r}$ una sucesión
de elementos de $\mathfrak{m}$, $\mathfrak{J}$ el ideal
$x_1A+\cdots+x_rA$. Las condiciones siguientes son equivalentes~:
\begin{enumerate}
  \item[\rm a)] $A$ es regular y los $x_i$ forman parte de un sistema
  regular de parámetros de $A$.
  \item[\rm a')] $A$ es regular y las clases $\overline{x}_i$ de los $x_i$
  módulo $\mathfrak{m}^2$ son linealmente independientes sobre $k$.
  \item[\rm b)] Los $x_i$ forman parte de un sistema de parámetros para $A$ y
  $A/\mathfrak{J}$ es un anillo regular.
\end{enumerate}
Además, cuando se verifican estas condiciones, $\mathfrak{J}$ es un ideal
primo.
\end{proposition}

\begin{proof}
La última aserción se deduce trivialmente de que $A/\mathfrak{J}$,
por ser regular, es íntegro (17.1.3). Pongamos $n=\dim(A)$. La equivalencia de a)
y a') es inmediata~; en efecto, las clases módulo $\mathfrak{m}^2$ de los elementos
de un sistema regular de parámetros forman una base de
$\mathfrak{m}/\mathfrak{m}^2$ sobre $k$, y recíprocamente, si los
$\overline{x}_i$ $(1\leq i\leq r)$ son linealmente independientes sobre $k$, se
pueden encontrar $n-r$ elementos $x_i\in\mathfrak{m}$ $(r+1\leq i\leq n)$ tales que
los $\overline{x}_i$ $(1\leq i\leq n)$ formen una base de
$\mathfrak{m}/\mathfrak{m}^2$, por lo que $(x_i)_{1\leq i\leq n}$ es un sistema
regular de parámetros. Para ver que a) implica b), observemos que, puesto que
los $x_i$ forman parte de un sistema de parámetros de $A$, se tiene
$\dim(A/\mathfrak{J})=n-r$ (16.3.6)~; sea
$\mathfrak{n}=\mathfrak{m}/\mathfrak{J}$ el ideal maximal de
$A/\mathfrak{J}$. Se tiene una sucesión exacta
\begin{equation}
\label{0.17.1.7.1-es}
\tag{17.1.7.1}
  0\longrightarrow(\mathfrak{m}^2+\mathfrak{J})/\mathfrak{m}^2
  \longrightarrow\mathfrak{m}/\mathfrak{m}^2
  \longrightarrow\mathfrak{n}/\mathfrak{n}^2\longrightarrow0
\end{equation}
puesto que $\mathfrak{n}^2=(\mathfrak{m}^2+\mathfrak{J})/\mathfrak{J}$~;
la hipótesis a) implica que
$\operatorname{rg}_k((\mathfrak{m}^2+\mathfrak{J})/\mathfrak{m}^2)=r$, por lo que
$\operatorname{rg}_k(\mathfrak{n}/\mathfrak{n}^2)=n-r
=\dim(A/\mathfrak{J})$, lo que demuestra b) en virtud de (17.1.1, b)).

Recíprocamente, para demostrar que b) implica a), observemos que el hecho de que los $x_i$
formen parte de un sistema de parámetros implica que
$\dim(A/\mathfrak{J})=n-r$ (16.3.6)~; la hipótesis de que $A/\mathfrak{J}$ es
regular implica, por otra parte,
$\operatorname{rg}_k(\mathfrak{n}/\mathfrak{n}^2)=n-r$ (17.1.1)~; además,
se tiene evidentemente
$\operatorname{rg}_k((\mathfrak{m}^2+\mathfrak{J})/\mathfrak{m}^2)\leq r$,
por lo que se deduce de (17.1.7.1) que
$\operatorname{rg}_k(\mathfrak{m}/\mathfrak{m}^2)\leq r+(n-r)=n$~; así, $A$
es regular en virtud de (16.2.6) y (17.1.1). Además, la relación
$\operatorname{rg}_k(\mathfrak{m}/\mathfrak{m}^2)=n$ implica entonces
$\operatorname{rg}_k((\mathfrak{m}^2+\mathfrak{J})/\mathfrak{m}^2)=r$, y por
consiguiente los $\overline{x}_i$ son linealmente independientes sobre $k$. C.Q.F.D.
\end{proof}

\oldpage[0\textsubscript{IV-1}]{42}

\begin{corollary}[17.1.8]
\label{0.17.1.8-es}
Sean $A$ un anillo local noetheriano, $\mathfrak{m}$ su ideal maximal, $t$ un
elemento de $\mathfrak{m}$. Las condiciones siguientes son equivalentes~:
\begin{enumerate}
  \item[\rm a)] $A/tA$ es regular y $t$ es no divisor de cero en $A$.
  \item[\rm b)] $A$ es regular y $t\notin\mathfrak{m}^2$.
\end{enumerate}
\end{corollary}

\begin{proof}
Es el caso particular $r=1$ de (17.1.7) (habida cuenta de que un elemento de
$\mathfrak{m}$ no divisor de cero forma parte de un sistema de parámetros
(16.4.1)).
\end{proof}

\begin{corollary}[17.1.9]
\label{0.17.1.9-es}
Sean $A$ un anillo local regular, $\mathfrak{m}$ su ideal maximal,
$\mathfrak{J}$ un ideal de $A$ contenido en $\mathfrak{m}$. Las condiciones
siguientes son equivalentes~:
\begin{enumerate}
  \item[\rm a)] El anillo $A/\mathfrak{J}$ es regular.
  \item[\rm b)] $\mathfrak{J}$ está engendrado por una sucesión
  $(x_i)_{1\leq i\leq r}$ que forma parte de un sistema regular de parámetros
  de $A$.
\end{enumerate}
\end{corollary}

\begin{proof}
Ya se ha visto (17.1.7) que b) implica a). Recíprocamente, supongamos que
$A/\mathfrak{J}$ es regular (lo que implica que $\mathfrak{J}$ es primo)
y sean $n=\dim(A)$, $n-r=\dim(A/\mathfrak{J})$; con la notación de
(17.1.7), la sucesión exacta (17.1.7.1) da
$\operatorname{rg}_k((\mathfrak{m}^2+\mathfrak{J})/\mathfrak{m}^2)=r$, luego
existe una sucesión $(x_i)_{1\leq i\leq r}$ de elementos de $\mathfrak{J}$ que
forma parte de un sistema regular de parámetros de $A$. Pongamos
$\mathfrak{J}'=x_1A+\cdots+x_rA$; de (17.1.7) resulta que $\mathfrak{J}'$
es un ideal primo de $A$ y que $A/\mathfrak{J}'$ es regular y de
dimensión $n-r$; pero como
$A/\mathfrak{J}=(A/\mathfrak{J}')/(\mathfrak{J}/\mathfrak{J}')$ y
$\mathfrak{J}/\mathfrak{J}'$ es primo en el anillo íntegro
$A/\mathfrak{J}'$, las dimensiones de $A/\mathfrak{J}$ y de
$A/\mathfrak{J}'$ solo pueden ser iguales si
$\mathfrak{J}=\mathfrak{J}'$ (16.1.2.2).
\end{proof}

\subsection{Recordatorios sobre la dimensión proyectiva y la dimensión inyectiva de los
módulos}
\label{subsection:0.17.2-es}

\begin{env}[17.2.1]
\label{0.17.2.1-es}
Sean $A$ un anillo y $M$ un $A$-módulo. Recordemos (M, VI, 2) que se denomina
\emph{dimensión proyectiva} (resp. \emph{dimensión inyectiva}) de $M$ y que
se denota $\operatorname{dim.proj}(M)$ o
$\operatorname{dim.proj}_A(M)$ (resp. $\operatorname{dim.inj}(M)$ o
$\operatorname{dim.inj}_A(M)$) al menor $n$ (entero o igual a $+\infty$)
tal que existe una resolución proyectiva por la izquierda (resp. una resolución
inyectiva por la derecha) de $M$ que sea \emph{de longitud $n$}. Es equivalente
(\emph{loc. cit.}) decir que $n$ es el menor número tal que para todo
$A$-módulo $N$, se tiene $\operatorname{Ext}^i_A(M,N)=0$ (resp.
$\operatorname{Ext}^i_A(N,M)=0$) para todo $i>n$, o tan solo para $i=n+1$.
Para todo \emph{factor directo} $M'$ de $M$, se tiene, pues,
$\operatorname{dim.proj}(M')\leq\operatorname{dim.proj}(M)$, ya que
$\operatorname{Ext}^i_A(M',N)$ es un factor directo de
$\operatorname{Ext}^i_A(M,N)$; asimismo,
$\operatorname{dim.inj}(M')\leq\operatorname{dim.inj}(M)$. Una condición
equivalente a $\operatorname{dim.proj}(M)=n$ (resp.
$\operatorname{dim.inj}(M)=n$) es que $n$ sea el menor número tal que,
para toda sucesión exacta
\[
  0\longrightarrow R\longrightarrow P_{n-1}\longrightarrow\cdots
  \longrightarrow P_0\longrightarrow M\longrightarrow0
\]
(resp.
$0\longrightarrow M\longrightarrow Q_0\longrightarrow\cdots
\longrightarrow Q_{n-1}\longrightarrow C\longrightarrow0$), donde todos los
$P_i$ para $0\leq i<n$ son proyectivos (resp. todos los $Q_i$ para $0\leq i<n$
son inyectivos), entonces $R$ es proyectivo (resp. $C$ es inyectivo).
\end{env}

\begin{remarks}[17.2.2]
\label{0.17.2.2-es}
\begin{enumerate}
  \item[\rm (i)] Decir que $M$ es un $A$-módulo \emph{proyectivo} (resp.
  \emph{inyectivo}) equivale, pues, a decir que
  $\operatorname{dim.proj}(M)=0$ (resp. $\operatorname{dim.inj}(M)=0$).
  \item[\rm (ii)] Sea $T$ un \emph{funtor covariante} aditivo de la categoría
  de los $A$-módulos en una categoría abeliana. De las definiciones de
  (17.2.1) y de la definición de los
  \oldpage[0\textsubscript{IV-1}]{43}
  funtores derivados resulta inmediatamente que si $\operatorname{dim.proj}(M)\leq n$ (resp.
  $\operatorname{dim.inj}(M)\leq n$), se tiene $L^iT(M)=0$ (resp.
  $R^iT(M)=0$) \emph{para todo $i>n$}.
  \item[\rm (iii)] Si se supone que $A$ es \emph{noetheriano} y $M$ es \emph{de tipo
  finito}, la última interpretación de la dimensión proyectiva dada en
  (17.2.1) muestra que, si $\operatorname{dim.proj}(M)=n$, $M$ admite una
  resolución por la izquierda de longitud $n$ formada por módulos proyectivos \emph{de tipo
  finito}. Si además $A$ es un anillo \emph{local} noetheriano, $M$ admite una
  resolución de longitud $n$ por módulos \emph{libres de tipo finito}, pues un
  $A$-módulo proyectivo de tipo finito es entonces libre (Bourbaki,
  \emph{Alg.~comm.}, cap.~II, \S~3, núm.~2, cor.~2 de la prop.~5).
\end{enumerate}
\end{remarks}

\begin{lemma}[17.2.3]
\label{0.17.2.3-es}
Sean $A$ un anillo y $M$ un $A$-módulo. Para que
$\operatorname{dim.inj}(M)\leq n$, es necesario y suficiente que para todo
$A$-módulo monógeno $N$, se tenga $\operatorname{Ext}^{n+1}_A(N,M)=0$.
\end{lemma}

\begin{proof}
Con la notación de (17.2.1), se trata de probar que $C$ es inyectivo;
para todo $A$-módulo $N$, $\operatorname{Ext}^{n+1}_A(N,M)$ es isomorfo a
$\operatorname{Ext}^1_A(N,C)$ (M, V, 7), luego se tiene
$\operatorname{Ext}^1_A(N,C)=0$ para todo $A$-módulo $N$ de la forma
$A/\mathfrak{J}$, donde $\mathfrak{J}$ es un ideal cualquiera de $A$. La sucesión
exacta de Ext muestra entonces que el homomorfismo canónico
$\operatorname{Hom}(A,C)\longrightarrow\operatorname{Hom}(\mathfrak{J},C)$ es
sobreyectivo para todo ideal $\mathfrak{J}$ de $A$, lo que implica que $C$ es un
$A$-módulo inyectivo (M, I, 3.2).
\end{proof}

\begin{lemma}[17.2.4]
\label{0.17.2.4-es}
Sean $A$ un anillo noetheriano y $M$ un $A$-módulo de tipo finito. Para que
$\operatorname{dim.proj}(M)\leq n$, es necesario y suficiente que para todo
$A$-módulo monógeno $N$, se tenga $\operatorname{Ext}^{n+1}_A(M,N)=0$.
\end{lemma}

\begin{proof}
En efecto, se sabe (M, VI, 2.5) que la condición
$\operatorname{dim.proj}(M)\leq n$ es equivalente a
$\operatorname{Ext}^{n+1}_A(M,N)=0$ para todo $A$-módulo $N$ \emph{de tipo
finito}; para ver que la condición del enunciado también es suficiente, se
razona por inducción sobre el número de generadores $m$ de $N$: existe un
submódulo $N_1$ de $N$ engendrado por $m-1$ elementos y tal que
$N_2=N/N_1$ es monógeno; de la sucesión exacta
$0\longrightarrow N_1\longrightarrow N\longrightarrow N_2\longrightarrow0$,
se deduce entonces la sucesión exacta
$\operatorname{Ext}^{n+1}_A(M,N_1)\longrightarrow
\operatorname{Ext}^{n+1}_A(M,N)\longrightarrow
\operatorname{Ext}^{n+1}_A(M,N_2)$, y la hipótesis de inducción muestra que la
condición del enunciado implica efectivamente $\operatorname{Ext}^{n+1}_A(M,N)=0$.
\end{proof}

\begin{corollary}[17.2.5]
\label{0.17.2.5-es}
Sean $A$ un anillo noetheriano y $M$ un $A$-módulo. Se tiene
\begin{equation}
\label{0.17.2.5.1-es}
\tag{17.2.5.1}
  \operatorname{dim.inj}_A(M)
  =\sup_{\mathfrak{m}}
  \bigl(\operatorname{dim.inj}_{A_{\mathfrak{m}}}(M_{\mathfrak{m}})\bigr)
\end{equation}
y, si $M$ es \emph{de tipo finito},
\begin{equation}
\label{0.17.2.5.2-es}
\tag{17.2.5.2}
  \operatorname{dim.proj}_A(M)
  =\sup_{\mathfrak{m}}
  \bigl(\operatorname{dim.proj}_{A_{\mathfrak{m}}}(M_{\mathfrak{m}})\bigr)
\end{equation}
donde $\mathfrak{m}$ recorre el conjunto de los ideales primos (o el conjunto de los
ideales maximales) de $A$.
\end{corollary}

\begin{proof}
En efecto, si $N$ es un $A$-módulo de tipo finito (por tanto, de presentación finita),
se tiene, para todo conjunto multiplicativo $S$ de $A$ y para todo $i\geq0$,
\[
  S^{-1}\operatorname{Ext}^i_A(N,M)
  =\operatorname{Ext}^i_{S^{-1}A}(S^{-1}N,S^{-1}M)
\]
por planitud, considerando una resolución libre de $M$ y utilizando el
hecho de que la relación precedente es válida para $i=0$ (Bourbaki,
\emph{Alg.~comm.}, cap.~II, \S~2, núm.~7, prop.~19). En particular,
$\operatorname{Ext}^i_{A_{\mathfrak{m}}}
(A_{\mathfrak{m}}/\mathfrak{J}A_{\mathfrak{m}},M_{\mathfrak{m}})
=(\operatorname{Ext}^i_A(A/\mathfrak{J},M))_{\mathfrak{m}}$ para todo ideal
primo $\mathfrak{m}$ y todo ideal $\mathfrak{J}$ de $A$; si se tienen
en cuenta (17.2.3) y el hecho de que todo
\oldpage[0\textsubscript{IV-1}]{44}
ideal de $A_{\mathfrak{m}}$ es de la forma $\mathfrak{J}A_{\mathfrak{m}}$
para un ideal adecuado $\mathfrak{J}$ de $A$, se deduce la fórmula
(17.2.5.1) de lo que precede y de Bourbaki, \emph{Alg.~comm.}, cap.~II,
\S~3, núm.~3, teorema~1. Se procede del mismo modo para (17.2.5.2), utilizando esta vez
(17.2.4) e intercambiando los papeles de $M$ y $N$.
\end{proof}

Para los anillos noetherianos y los módulos de tipo finito, el estudio de la
dimensión proyectiva o de la dimensión inyectiva queda, pues, reducido por
(17.2.5) al caso de los anillos \emph{locales}. Se tiene entonces el

\begin{lemma}[17.2.6]
\label{0.17.2.6-es}
Sean $A$ un anillo local noetheriano, $k$ su cuerpo residual y $M$ un
$A$-módulo de tipo finito. Para que $\operatorname{dim.proj}(M)\leq n$, es
necesario que $\operatorname{Tor}^A_i(M,k)=0$ para $i>n$, y basta que
$\operatorname{Tor}^A_{n+1}(M,k)=0$.
\end{lemma}

\begin{proof}
La necesidad es un caso particular de la observación (17.2.2 (ii)), aplicada al
funtor covariante $M\mapsto k\otimes_A M$. Para probar que la condición es
suficiente, se trata, con la notación de (17.2.1), de establecer que $R$ es
proyectivo cuando se supone que los $P_i$ son de tipo finito; ahora bien,
$\operatorname{Tor}^A_{n+1}(M,k)$ es isomorfo a
$\operatorname{Tor}^A_1(R,k)$ (M, V, 7); y se sabe que, puesto que $R$ es de
tipo finito, la condición $\operatorname{Tor}^A_1(R,k)=0$ implica que $R$ es
libre (Bourbaki, \emph{Alg.~comm.}, cap.~II, \S~3, núm.~2, cor.~2 de la
prop.~5).
\end{proof}

\begin{corollary}[17.2.7]
\label{0.17.2.7-es}
Bajo las hipótesis de (17.2.6), sea $x$ un elemento $M$-regular de $A$
perteneciente al ideal maximal de $A$; entonces se tiene
\begin{equation}
\label{0.17.2.7.1-es}
\tag{17.2.7.1}
  \operatorname{dim.proj}(M/xM)=\operatorname{dim.proj}(M)+1.
\end{equation}
\end{corollary}

\begin{proof}
En efecto, de la sucesión exacta
$0\longrightarrow M\xrightarrow{x}M\longrightarrow M/xM\longrightarrow0$
definida por la multiplicación por $x$ en $M$, se deduce la sucesión exacta de
Tor:
\[
  \operatorname{Tor}^A_i(M,k)\xrightarrow{x}
  \operatorname{Tor}^A_i(M,k)\longrightarrow
  \operatorname{Tor}^A_i(M/xM,k)\longrightarrow
  \operatorname{Tor}^A_{i-1}(M,k)\xrightarrow{x}
  \operatorname{Tor}^A_{i-1}(M,k)
\]
para todo $i\geq1$. Ahora bien, para todo $A$-módulo $N$, la homotecia de razón $x$
en $M\otimes_A N$ procede tanto de la homotecia de razón $x$ en
$M$ como de la homotecia de razón $x$ en $N$; se concluye inmediatamente,
puesto que la homotecia de razón $x$ en $k$ es nula por definición, que,
para todo $i$, el homomorfismo
$\operatorname{Tor}^A_i(M,k)\xrightarrow{x}\operatorname{Tor}^A_i(M,k)$ es
\emph{nulo}; dicho de otro modo, se tiene la sucesión exacta
\[
  0\longrightarrow\operatorname{Tor}^A_i(M,k)\longrightarrow
  \operatorname{Tor}^A_i(M/xM,k)\longrightarrow
  \operatorname{Tor}^A_{i-1}(M,k)\longrightarrow0.
\]
Si $\operatorname{dim.proj}(M)=n$, se tiene
$\operatorname{Tor}^A_n(M,k)\neq0$ y
$\operatorname{Tor}^A_{n+1}(M,k)=\operatorname{Tor}^A_{n+2}(M,k)=0$ en virtud
de (17.2.6). De lo que precede resulta, pues, que se tiene
$\operatorname{Tor}^A_{n+1}(M/xM,k)\neq0$ y
$\operatorname{Tor}^A_{n+2}(M/xM,k)=0$, luego
$\operatorname{dim.proj}(M/xM)=n+1$ por (17.2.6).
\end{proof}

\begin{proposition}[17.2.8]
\label{0.17.2.8-es}
\textnormal{(M.~Auslander)} Sea $A$ un anillo. Las condiciones siguientes son
equivalentes:
\begin{enumerate}
  \item[\rm a)] Todo $A$-módulo es de dimensión proyectiva $\leq n$.
  \item[\rm a')] Todo $A$-módulo de tipo finito es de dimensión proyectiva
  $\leq n$.
  \item[\rm b)] Todo $A$-módulo es de dimensión inyectiva $\leq n$.
  \item[\rm c)] Para todo par de $A$-módulos $M$, $N$, se tiene
  $\operatorname{Ext}^{n+1}_A(M,N)=0$.
  \item[\rm c')] Para todo par de $A$-módulos $M$, $N$ tal que $M$ sea de
  tipo finito (o solamente monógeno), se tiene
  $\operatorname{Ext}^{n+1}_A(M,N)=0$.
\end{enumerate}
\end{proposition}

\begin{proof}
\oldpage[0\textsubscript{IV-1}]{45}
Esto resulta inmediatamente de (17.2.1) y (17.2.3).
\end{proof}

El menor número $n$ (entero o $+\infty$) para el cual se satisfacen las condiciones
equivalentes de (17.2.8) se denomina \emph{dimensión cohomológica
global} (o simplemente \emph{dimensión cohomológica}) de $A$
y se denota por $\operatorname{dim.coh}(A)$.

\begin{proposition}[17.2.9]
\label{0.17.2.9-es}
Sea $A$ un anillo noetheriano. Las condiciones siguientes son equivalentes~:
\begin{enumerate}
  \item[\rm a)] $\operatorname{dim.coh}(A)\leq n$.
  \item[\rm b)] Todo $A$-módulo de tipo finito es de dimensión inyectiva
  $\leq n$.
  \item[\rm c)] Para todo par de $A$-módulos de tipo finito, $M$, $N$, se tiene
  $\operatorname{Ext}^{n+1}_A(M,N)=0$.
\end{enumerate}
\end{proposition}

\begin{proof}
Esto resulta inmediatamente de la definición y de (17.2.4).
\end{proof}

\begin{corollary}[17.2.10]
\label{0.17.2.10-es}
Si $A$ es un anillo noetheriano, se tiene
\begin{equation}
\label{0.17.2.10.1-es}
\tag{17.2.10.1}
  \operatorname{dim.coh}(A)
  =\sup_{\mathfrak{m}}
  \bigl(\operatorname{dim.coh}(A_{\mathfrak{m}})\bigr)
\end{equation}
donde $\mathfrak{m}$ recorre el espectro de $A$ (o el conjunto de los ideales
maximales de $A$).
\end{corollary}

\begin{proof}
Esto se deduce de (17.2.9) y (17.2.5).
\end{proof}

\begin{proposition}[17.2.11]
\label{0.17.2.11-es}
Sean $A$ un anillo local noetheriano y $k$ su cuerpo residual. Para que
$\operatorname{dim.coh}(A)\leq n$, es necesario que
$\operatorname{Tor}^A_i(k,k)=0$, para $i>n$, y basta con que
$\operatorname{Tor}^A_{n+1}(k,k)=0$.
\end{proposition}

\begin{proof}
Habida cuenta de (17.2.6), basta demostrar que las relaciones
$\operatorname{dim.coh}(A)\leq n$ y
$\operatorname{dim.proj}_A(k)\leq n$ son equivalentes. Es claro que la
primera implica la segunda por definición. Recíprocamente, si
$\operatorname{dim.proj}_A(k)\leq n$, se tiene
$\operatorname{Tor}^A_i(M,k)=0$ para $i>n$ y para todo $A$-módulo $M$ en
virtud de (17.2.2, (ii))~; por tanto, $\operatorname{dim.proj}(M)\leq n$, lo que
demuestra la proposición en virtud de (17.2.8).
\end{proof}

\begin{corollary}[17.2.12]
\label{0.17.2.12-es}
Sea $A$ un anillo noetheriano. Para que se tenga
$\operatorname{dim.coh}(A)\leq n$, es necesario que, para todo ideal
maximal $\mathfrak{m}$ de $A$, se tenga
$\operatorname{Tor}^{A_{\mathfrak{m}}}_i(A/\mathfrak{m},A/\mathfrak{m})=0$
para $i>n$, y basta con que estas relaciones se verifiquen para $i=n+1$.
\end{corollary}

\begin{proof}
Esto resulta inmediatamente de (17.2.11) y (17.2.10).
\end{proof}

\begin{proposition}[17.2.13]
\label{0.17.2.13-es}
Sean $A$, $B$ dos anillos locales noetherianos y
$\varphi:A\longrightarrow B$ un homomorfismo local que hace de $B$ un
$A$-módulo plano. Entonces se tiene
\begin{equation}
\label{0.17.2.13.1-es}
\tag{17.2.13.1}
  \operatorname{dim.coh}(A)\leq\operatorname{dim.coh}(B).
\end{equation}
\end{proposition}

\begin{proof}
Supongamos que $\operatorname{dim.coh}(B)=n$ es finita~; basta demostrar
que, para todo par $(M,N)$ de $A$-módulos de tipo finito, se tiene
$\operatorname{Tor}^A_i(M,N)=0$ para $i>n$ (17.2.11). Como $B$ es un
$A$-módulo fielmente plano ($0_{\rm I}$, 6.6.2), demostrar esto equivale
($0_{\rm I}$, 6.4.1) a demostrar que se tiene
\begin{equation}
\label{0.17.2.13.2-es}
\tag{17.2.13.2}
  \operatorname{Tor}^A_i(M,N)\otimes_A B=0.
\end{equation}
Ahora bien, si $L_\bullet=(L_j)$ es una resolución derecha de $M$ por $A$-módulos
libres, del hecho de que $B$ sea un $A$-módulo plano se deduce que
$L_\bullet\otimes_A B=(L_j\otimes_A B)$ es una resolución derecha del
$B$-módulo $M\otimes_A B$ por $B$-módulos libres~; además, se tiene
$(L_j\otimes_A B)\otimes_B(N\otimes_A B)
=(L_j\otimes_A N)\otimes_A B$, de donde se concluye inmediatamente que el primer miembro
\oldpage[0\textsubscript{IV-1}]{46}
de (17.2.13.2) es igual a
$\operatorname{Tor}^B_i(M\otimes_A B,N\otimes_A B)$~; la hipótesis sobre $B$
implica que este $B$-módulo es nulo para $i>n$ (17.2.2, (ii)), de donde resulta la
conclusión.
\end{proof}

\begin{env}[17.2.14]
\label{0.17.2.14-es}
Sean $(X,\mathscr{O}_X)$ un espacio anillado y $\mathscr{F}$ un
$\mathscr{O}_X$-Módulo~; se llama \emph{dimensión proyectiva} (resp.
\emph{inyectiva}) \emph{puntual} de $\mathscr{F}$ y se denota por
$\operatorname{dim.proj}(\mathscr{F})$ (resp.
$\operatorname{dim.inj}(\mathscr{F})$) el número
$\sup_{x\in X}(\operatorname{dim.proj}(\mathscr{F}_x))$ (resp.
$\sup_{x\in X}(\operatorname{dim.inj}(\mathscr{F}_x))$. Se llama
\emph{dimensión cohomológica puntual de $X$} y se denota por
$\operatorname{dim.coh}(X)$ el número
$\sup_{x\in X}(\operatorname{dim.coh}(\mathscr{O}_x))$. Se deduce de
(17.2.5) y (17.2.10) que, cuando $A$ es un anillo noetheriano, $M$ un
$A$-módulo de tipo finito y $X=\operatorname{Spec}(A)$, se tiene
$\operatorname{dim.proj}(\widetilde{M})=\operatorname{dim.proj}(M)$,
$\operatorname{dim.inj}(\widetilde{M})=\operatorname{dim.inj}(M)$ y
$\operatorname{dim.coh}(X)=\operatorname{dim.coh}(A)$. Se llama
\emph{dimensión proyectiva} (resp. \emph{inyectiva}) \emph{de $\mathscr{F}$ en
un punto $x\in X$} a la dimensión proyectiva (resp. inyectiva) de
$\mathscr{F}_x$, y \emph{dimensión cohomológica de $X$ en un punto $x$} a la
dimensión cohomológica de $\mathscr{O}_x$.
\end{env}

\begin{proposition}[17.2.15]
\label{0.17.2.15-es}
Sean $X$, $Y$ dos espacios anillados en anillos locales noetherianos y
$f:X\longrightarrow Y$ un morfismo plano. Si
$\operatorname{dim.coh}(X)\leq n$, entonces $Y$ es de dimensión cohomológica
$\leq n$ en todo punto de $f(X)$.
\end{proposition}

\begin{proof}
Esto resulta inmediatamente de (17.2.13).
\end{proof}

\subsection{Teoría cohomológica de los anillos regulares}
\label{subsection:0.17.3-es}

\begin{theorem}[17.3.1]
\label{0.17.3.1-es}
\textnormal{(Hilbert-Serre)} Sea $A$ un anillo local noetheriano. Para que $A$
sea de dimensión cohomológica finita, es necesario y suficiente que $A$ sea
regular~; se tiene entonces
\begin{equation}
\label{0.17.3.1.1-es}
\tag{17.3.1.1}
  \operatorname{dim.coh}(A)=\dim(A).
\end{equation}
\end{theorem}

\begin{proof}
Supongamos que $A$ es regular~; sean $\mathfrak{m}$ su ideal maximal y
$\mathbf{x}=(x_i)_{1\leq i\leq n}$ un sistema regular de parámetros de $A$
(17.1.6)~; consideremos el complejo del álgebra exterior
$K_\bullet(\mathbf{x})$ (\textbf{III}, 1.1.1), que está formado por $A$-módulos
libres, con $K_i(\mathbf{x})=0$ para $i>n$~; como $(x_i)$ es una sucesión
regular, se tiene $H_i(K_\bullet(\mathbf{x}))=0$ para $i>0$ (\textbf{III},
1.1.4 y 1.1.3.3) y
$H_0(K_\bullet(\mathbf{x}))=A/(x_1A+\cdots+x_nA)=A/\mathfrak{m}=k$~; los
$K_i(\mathbf{x})$ forman, por tanto, una \emph{resolución libre} de $k$ de
longitud $n$. Ahora bien, el hecho de que los $x_i$ pertenezcan a $\mathfrak{m}$
implica inmediatamente que, en el complejo
$K_\bullet(\mathbf{x},k)=K_\bullet(\mathbf{x})\otimes_A k$, el operador de borde
es nulo en todos los grados, de modo que se tiene, por definición,
$\operatorname{Tor}^A_i(k,k)=\bigwedge^i(k^n)$~; la igualdad (17.3.1.1) resulta
entonces inmediatamente de (17.2.11) (este resultado es esencialmente el «~teorema de las
sizigias~» de Hilbert).

Demostremos ahora que, si $A$ es un anillo local noetheriano,
$\mathfrak{m}$ su ideal maximal, y si
$\operatorname{dim.proj}_A(\mathfrak{m})$ es \emph{finita}, entonces $A$ es
regular, lo que terminará la demostración de (17.3.1). Procederemos por
inducción sobre $n=\operatorname{rg}_k(\mathfrak{m}/\mathfrak{m}^2)$. Para
$n=0$, se tiene $\mathfrak{m}=0$ y la aserción es trivial.
\end{proof}

\begin{lemma}[17.3.1.2]
\label{0.17.3.1.2-es}
\textnormal{(Nagata)} Si todo elemento de
$\mathfrak{m}-\mathfrak{m}^2$ es divisor de cero en $A$, existe
$c\neq0$ en $A$ tal que $c\mathfrak{m}=0$ (dicho de otro modo, se tiene
$\mathfrak{m}\in\operatorname{Ass}(A)$).
\end{lemma}

\begin{proof}
Podemos limitarnos al caso en que $\mathfrak{m}\neq0$, y por tanto
$\mathfrak{m}\neq\mathfrak{m}^2$. La hipótesis implica que
$\mathfrak{m}-\mathfrak{m}^2$ está contenido en la unión de los ideales
$\mathfrak{p}_i$ de $\operatorname{Ass}(A)$ (Bourbaki,
\emph{Alg.~comm.}, cap.~IV, \S~1, no~1, cor.~3 de la prop.~2)~; por tanto,
$\mathfrak{m}$ está contenido en la unión de $\mathfrak{m}^2$
\oldpage[0\textsubscript{IV-1}]{47}
y de los $\mathfrak{p}_i$, y por tanto en uno de los $\mathfrak{p}_i$, puesto que
$\mathfrak{m}\not\subset\mathfrak{m}^2$ (Bourbaki, \emph{Alg.~comm.},
cap.~II, \S~1, no~1, prop.~2)~; como $\mathfrak{m}$ es maximal, esto demuestra el
lema.
\end{proof}

\begin{lemma}[17.3.1.3]
\label{0.17.3.1.3-es}
Si $a\in\mathfrak{m}-\mathfrak{m}^2$, $\mathfrak{m}/Aa$ es isomorfo a un
factor directo de $\mathfrak{m}/a\mathfrak{m}$.
\end{lemma}

\begin{proof}
Existe un sistema minimal de generadores de $\mathfrak{m}$ que contiene a $a$
(cuyas imágenes en $\mathfrak{m}/\mathfrak{m}^2$ forman una base de este
$k$-espacio vectorial)~; sea $\mathfrak{b}$ el ideal engendrado por los elementos
de este sistema distintos de $a$. Como la relación $xa\in\mathfrak{b}$ implica
$x\in\mathfrak{m}$ (considerando la imagen de $xa$ en
$\mathfrak{m}/\mathfrak{m}^2$), se tiene
$\mathfrak{b}\cap Aa\subset a\mathfrak{m}$, de donde se obtiene un homomorfismo
$\mathfrak{b}/(\mathfrak{b}\cap Aa)\longrightarrow
\mathfrak{m}/a\mathfrak{m}$ deducido de la inyección
$\mathfrak{b}\longrightarrow\mathfrak{m}$ por paso a los cocientes, y que
es \emph{inyectivo}. Además, $\mathfrak{b}+Aa=\mathfrak{m}$, y
el homomorfismo compuesto
\[
  \mathfrak{m}/Aa=(\mathfrak{b}+Aa)/Aa
  \xrightarrow{\sim}\mathfrak{b}/(\mathfrak{b}\cap Aa)
  \longrightarrow\mathfrak{m}/a\mathfrak{m}
  \longrightarrow\mathfrak{m}/Aa
\]
es la identidad~; de donde resulta el lema.
\end{proof}

\begin{lemma}[17.3.1.4]
\label{0.17.3.1.4-es}
Sean $A$ un anillo local noetheriano, $\mathfrak{m}$ su ideal maximal, $E$
un $A$-módulo de tipo finito y de dimensión proyectiva finita. Si
$a\in\mathfrak{m}$ es $A$-regular y $E$-regular, entonces $E/aE$ es un
$(A/aA)$-módulo de dimensión proyectiva finita, a lo sumo igual a
$\operatorname{dim.proj}_A(E)$.
\end{lemma}
\begin{proof}
Razonemos por inducción sobre $h=\operatorname{dim.proj}_A(E)$, siendo trivial
el caso $h=0$, puesto que entonces $E$ es un $A$-módulo proyectivo, y por tanto
$E/aE$ es un $(A/aA)$-módulo proyectivo. Existe una sucesión exacta
\[
  0\longrightarrow N\longrightarrow L\longrightarrow E\longrightarrow0
\]
donde $L$ es libre y $\operatorname{dim.proj}_A(N)=h-1$ (17.2.2, (iii)), siendo
$N$ de tipo finito. Además, la sucesión
\[
  0\longrightarrow N/aN\longrightarrow L/aL\longrightarrow E/aE
  \longrightarrow0
\]
es exacta (15.1.18). Como $a$ es $A$-regular, también es $N$-regular, puesto
que $L$ es libre; la hipótesis de inducción implica que $N/aN$ es un
$(A/aA)$-módulo de dimensión proyectiva $\leq h-1$, y como $L/aL$ es un
$(A/aA)$-módulo libre, $E/aE$ es un $(A/aA)$-módulo de dimensión proyectiva
$\leq h$.
\end{proof}

Examinemos ahora dos casos:

\medskip\noindent\textbf{I.} --- Supongamos primero que todo elemento de
$\mathfrak{m}-\mathfrak{m}^2$ es divisor de cero en $A$, en cuyo caso
(17.3.1.2) existe $c\neq0$ en $A$ tal que $c\mathfrak{m}=0$. Demostremos
entonces que $\mathfrak{m}=0$. Si no fuera así, observemos primero que
$\mathfrak{m}$ no podría ser un $A$-módulo proyectivo, pues sería libre
(Bourbaki, \emph{Alg.~comm.}, cap.~II, \S~3, n.º~2, cor.~2 de la prop.~5),
lo que contradice la relación $c\mathfrak{m}=0$. Se tendría, por tanto,
$n=\operatorname{dim.coh}(A)\geq1$. Como
$\mathfrak{m}\in\operatorname{Ass}(A)$, existiría una sucesión exacta de
$A$-homomorfismos
\[
  0\longrightarrow k\longrightarrow A\longrightarrow E\longrightarrow0.
\]
Pero esto es absurdo, pues, en virtud de la relación $n\geq1$, la sucesión
exacta de los Tor daría la sucesión exacta
$0\longrightarrow\operatorname{Tor}^A_{n+1}(E,k)\longrightarrow
\operatorname{Tor}^A_n(k,k)\longrightarrow0$; ahora bien, se tiene
$\operatorname{Tor}^A_{n+1}(E,k)=0$ (17.2.2, (ii)) y
$\operatorname{Tor}^A_n(k,k)\neq0$ (17.2.11), y hemos llegado a una
contradicción.

\medskip\noindent\textbf{II.} --- Podemos, pues, limitarnos al caso en que
existe $a\in\mathfrak{m}-\mathfrak{m}^2$ que es un elemento $A$-regular y, por
tanto, también $\mathfrak{m}$-regular. Consideremos el anillo $A'=A/aA$ y su
ideal
\oldpage[0\textsubscript{IV-1}]{48}
máximo $\mathfrak{m}'=\mathfrak{m}/aA$; es claro que
$\operatorname{rg}_k(\mathfrak{m}'/\mathfrak{m}'^2)=n-1$. En virtud de
(17.3.1.4), $\mathfrak{m}/a\mathfrak{m}$ es un $A'$-módulo de dimensión
proyectiva finita, luego lo mismo ocurre con
$\mathfrak{m}'=\mathfrak{m}/aA$, que es un sumando directo suyo (17.3.1.3 y
17.2.1). La hipótesis de inducción implica, por tanto, que $A'=A/aA$ es
regular, lo que demuestra, por (17.1.8), que $A$ es regular.

\begin{corollary}[17.3.2]
\label{0.17.3.2-es}
Si $A$ es un anillo local regular, $A_{\mathfrak{p}}$ es regular para todo
ideal primo $\mathfrak{p}$ de $A$.
\end{corollary}

\begin{proof}
En efecto, se ha visto (17.2.10) que
$\operatorname{dim.coh}(A_{\mathfrak{p}})\leq
\operatorname{dim.coh}(A)$, luego la conclusión resulta inmediatamente de
(17.3.1).
\end{proof}

\begin{proposition}[17.3.3]
\label{0.17.3.3-es}
Sean $A$, $B$ dos anillos locales noetherianos,
$\varphi:A\longrightarrow B$ un homomorfismo local, $\mathfrak{m}$ (resp.
$\mathfrak{n}$) el ideal máximo de $A$ (resp. $B$),
$k=A/\mathfrak{m}$ (resp. $k'=B/\mathfrak{n}$) el cuerpo residual de $A$
(resp. $B$). Se tiene, por tanto, un homomorfismo canónico de $k'$-espacios
vectoriales
\begin{equation}
\label{0.17.3.3.1-es}
\tag{17.3.3.1}
  \psi:(\mathfrak{m}/\mathfrak{m}^2)\otimes_k k'
  \longrightarrow\mathfrak{n}/\mathfrak{n}^2.
\end{equation}
\begin{enumerate}
  \item[\rm (i)] Si $B$ es regular y $B$ es un $A$-módulo plano, $A$ es
  regular.
  \item[\rm (ii)] Las condiciones siguientes son equivalentes:
  \begin{enumerate}
    \item[\rm a)] $B$ es regular y el homomorfismo $\psi$ (17.3.3.1) es
    inyectivo.
    \item[\rm b)] $A$ y $B$ son regulares, y para un sistema regular de
    parámetros $(x_i)_{1\leq i\leq m}$ de $A$, los $\varphi(x_i)$ forman parte
    de un sistema regular de parámetros de $B$ (en cuyo caso esta propiedad se
    verifica para \emph{todo} sistema regular de parámetros de $A$).
    \item[\rm c)] $B$ y $B\otimes_A k$ son regulares, y $B$ es un $A$-módulo
    plano.
    \item[\rm d)] $A$ y $B\otimes_A k$ son regulares, y $B$ es un $A$-módulo
    plano.
    \item[\rm e)] $A$ y $B\otimes_A k$ son regulares, y se tiene
    \begin{equation}
    \label{0.17.3.3.2-es}
    \tag{17.3.3.2}
      \dim(B)=\dim(A)+\dim(B\otimes_A k).
    \end{equation}
  \end{enumerate}
\end{enumerate}
\end{proposition}

\begin{proof}
\textnormal{(i)} Se tiene
$\operatorname{dim.coh}(A)\leq\operatorname{dim.coh}(B)$ por (17.2.13), luego
basta aplicar (17.3.1).

\textnormal{(ii)} Cuando se supone que $A$ es regular y $(x_i)$ es un sistema
regular de parámetros de $A$, decir que $B$ es un $A$-módulo \emph{plano}
equivale, en virtud de (15.1.21), a decir que la sucesión de los
$y_i=\varphi(x_i)$ es $B$-regular (puesto que $B\otimes_A k$ es un
$k$-módulo plano). Por otra parte, como $\dim(A)=m$ y la sucesión $(y_i)$
engendra el ideal $\mathfrak{m}B$, la relación (17.3.3.2), que también se
escribe $\dim(B/\mathfrak{m}B)=\dim(B)-m$, equivale a decir que la sucesión
$(y_i)_{1\leq i\leq m}$ forma parte de un sistema de parámetros de $B$
(16.3.7). Se ve, por tanto, que las condiciones b), d) y e) equivalen,
respectivamente, a las siguientes:
\begin{enumerate}
  \item[\rm b')] $A$ y $B$ son regulares y $(y_i)_{1\leq i\leq m}$ forma
  parte de un sistema regular de parámetros de $B$.
  \item[\rm d')] $A$ es regular,
  $B/(\sum_{i=1}^m y_iB)$ es regular y la sucesión $(y_i)$ es
  $B$-regular.
  \item[\rm e')] $A$ es regular,
  $B/(\sum_{i=1}^m y_iB)$ es regular y la sucesión $(y_i)$ forma parte de un
  sistema de parámetros de $B$.
\end{enumerate}
Ahora bien, b') y e') son equivalentes en virtud de (17.1.7), y como d')
implica e') (16.4.1) y se sigue de b') (17.1.7), es equivalente a ellas.
La conjunción
\oldpage[0\textsubscript{IV-1}]{49}
de b) y d) implica trivialmente c), y, en virtud de (i), c) implica d); se ha
probado, por tanto, la equivalencia de b), c), d) y e). Además, es claro que
b) implica a), pues las clases de los $x_i$ en
$\mathfrak{m}/\mathfrak{m}^2$ forman entonces una base de este $k$-espacio
vectorial y las clases de los $y_i$ en $\mathfrak{n}/\mathfrak{n}^2$ un
sistema libre de este $k'$-espacio vectorial. Queda, pues, probar que a)
implica e). Pongamos $V=\mathfrak{m}/\mathfrak{m}^2$,
$W=\mathfrak{n}/\mathfrak{n}^2$; es inmediato que se tiene
$\psi(V\otimes_k k')=(\mathfrak{n}^2+\mathfrak{m}B)/\mathfrak{n}^2$. Por una
parte, en virtud de (16.3.9),
\begin{equation}
\label{0.17.3.3.3-es}
\tag{17.3.3.3}
  \dim(B)\leq\dim(A)+\dim(B\otimes_A k)~;
\end{equation}
en segundo lugar, según (16.2.6), se tiene
\begin{equation}
\label{0.17.3.3.4-es}
\tag{17.3.3.4}
  \dim(A)\leq\operatorname{rg}_{k'}(V\otimes_k k').
\end{equation}
Por último, en el anillo local $B\otimes_A k=B/\mathfrak{m}B$,
$\mathfrak{n}'=\mathfrak{n}/\mathfrak{m}B$ es el ideal máximo, $k'$ el cuerpo
residual y $\mathfrak{n}'/\mathfrak{n}'^2$ es isomorfo a
$\mathfrak{n}/(\mathfrak{n}^2+\mathfrak{m}B)
=W/\psi(V\otimes_k k')$; se tiene, por tanto, según (16.2.6),
\begin{equation}
\label{0.17.3.3.5-es}
\tag{17.3.3.5}
  \dim(B\otimes_A k)
  \leq\operatorname{rg}_{k'}W
  -\operatorname{rg}_{k'}\psi(V\otimes_k k').
\end{equation}
Finalmente, puesto que se supone que $B$ es regular, se tiene
$\dim(B)=\operatorname{rg}_{k'}W$ (17.1.1); se concluye, por tanto, de
(17.3.3.3), (17.3.3.4) y (17.3.3.5) que se tiene
\[
  \operatorname{rg}_{k'}W
  \leq\dim(A)+\dim(B\otimes_A k)
  \leq\operatorname{rg}_{k'}W
  +\operatorname{rg}_{k'}(V\otimes_k k')
  -\operatorname{rg}_{k'}\psi(V\otimes_k k')
\]
y decir que los dos términos extremos de esta desigualdad son iguales
significa que $\psi$ es \emph{inyectivo}. La condición a) implica, por tanto,
necesariamente que en cada una de las relaciones (17.3.3.3), (17.3.3.4) y
(17.3.3.5), los dos miembros son \emph{iguales}; ahora bien, la igualdad en
(17.3.3.4) (resp. (17.3.3.5)) significa que $A$ (resp. $B\otimes_A k$) es
\emph{regular} (17.1.1); queda así demostrado que a) implica e).
\end{proof}

\begin{proposition}[17.3.4]
\label{0.17.3.4-es}
Sean $A$ un anillo local regular de dimensión $n$, $\mathfrak{m}$ su ideal
máximo. Para todo $A$-módulo no nulo de tipo finito $M$, se tiene
\begin{equation}
\label{0.17.3.4.1-es}
\tag{17.3.4.1}
  \operatorname{prof}(M)+\operatorname{dim.proj}(M)=n.
\end{equation}
\end{proposition}

\begin{proof}
Razonemos por inducción sobre $m=\operatorname{prof}(M)$. Si $m=0$, se sabe
(16.4.6, (i)) que existe un sub-$A$-módulo $N$ de $M$ isomorfo a
$k=A/\mathfrak{m}$; aplicando la sucesión exacta de los Tor a la sucesión
exacta $0\longrightarrow N\longrightarrow M\longrightarrow M/N\longrightarrow0$,
se obtiene una sucesión exacta
\[
  \operatorname{Tor}^A_{n+1}(M/N,k)\longrightarrow
  \operatorname{Tor}^A_n(N,k)\longrightarrow
  \operatorname{Tor}^A_n(M,k),
\]
y la hipótesis de que $A$ es regular implica
$\operatorname{Tor}^A_{n+1}(M/N,k)=0$ ((17.3.1.1) y (17.2.6)), luego
$\operatorname{Tor}^A_n(k,k)$ es isomorfo a un sub-$A$-módulo de
$\operatorname{Tor}^A_n(M,k)$; aplicando de nuevo (17.3.1.1) y (17.2.6), se
ve que $\operatorname{Tor}^A_n(M,k)\neq0$, luego
$\operatorname{dim.proj}(M)\geq n$ (17.2.6); pero como
$\operatorname{dim.proj}(M)\leq n$ por (17.3.1.1), se tiene efectivamente
$\operatorname{dim.proj}(M)=n$.

Supongamos ahora que $m>0$, y sea $x$ un elemento $M$-regular perteneciente a
$\mathfrak{m}$; se sabe entonces (16.4.6, (i)) que se tiene
$\operatorname{prof}(M/xM)=m-1$, y, por otra parte, (17.2.7),
$\operatorname{dim.proj}(M/xM)=\operatorname{dim.proj}(M)+1$; la hipótesis de
inducción prueba inmediatamente la relación (17.3.4.1).
\end{proof}

\oldpage[0\textsubscript{IV-1}]{50}

\begin{corollary}[17.3.5]
\label{0.17.3.5-es}
\begin{enumerate}
  \item[\rm(i)] Sea $A$ un anillo local regular de dimensión $n$; para que un
  $A$-módulo $M\neq0$ de tipo finito sea libre, es necesario y suficiente que
  $M$ sea un $A$-módulo de Cohen-Macaulay de dimensión $n$.
  \item[\rm(ii)] Sean $A$ un anillo local regular, $B$ un anillo local,
  $\rho:A\to B$ un homomorfismo local que hace de $B$ un $A$-módulo de tipo
  finito. Para que $B$ sea un $A$-módulo libre, es necesario y suficiente que
  $B$ sea un anillo de Cohen-Macaulay y que $\rho$ sea inyectivo (o,
  equivalentemente (16.1.5), que $\dim(B)=\dim(A)$).
\end{enumerate}
\end{corollary}

\begin{proof}
(i) Esto resulta de (17.3.4.1) y del hecho de que, para un $A$-módulo de tipo
finito, es lo mismo decir que este módulo es proyectivo que decir que es libre
(Bourbaki, \emph{Alg. comm.}, cap.~II, \S~3, n\textsuperscript{o}~2, cor.~2 de
la prop.~5); los $A$-módulos libres de tipo finito $M$ se caracterizan, por
tanto, por la relación $\operatorname{dim.proj}(M)=0$ (17.2.2, (i)).

(ii) Decir que $B$ es un anillo de Cohen-Macaulay equivale a decir que $B$ es
un $A$-módulo de Cohen-Macaulay (16.5.3), luego basta aplicar (i), puesto que
$\dim(B)=\dim(A)$ (16.1.5).
\end{proof}

\begin{env}[17.3.6]
\label{0.17.3.6-es}
Se dice que un anillo noetheriano $A$ es \emph{regular} si, para todo ideal
primo $\mathfrak{p}$ de $A$, el anillo local $A_{\mathfrak{p}}$ es regular;
cuando el propio $A$ es un anillo local, se sigue de (17.3.2) que esta
definición es equivalente a la de (17.1.2). Para que $A$ sea regular, basta,
en virtud de (17.3.2), que $A_{\mathfrak{m}}$ sea regular para todo ideal
máximo $\mathfrak{m}$ de $A$. Además, se sigue inmediatamente de esta
definición que, para todo conjunto multiplicativamente cerrado $S$ de $A$,
$S^{-1}A$ es regular.
\end{env}

\begin{proposition}[17.3.7]
\label{0.17.3.7-es}
Si $A$ es un anillo noetheriano regular, todo anillo de polinomios
$A[T_1,\ldots,T_n]$ es regular.
\end{proposition}

\begin{proof}
Es evidentemente suficiente probar que el anillo de polinomios $B=A[T]$ es
regular; como $B$ es un $A$-módulo libre, para todo ideal primo
$\mathfrak{q}$ de $B$, $B_{\mathfrak{q}}$ es un
$A_{\mathfrak{p}}$-módulo plano, donde $\mathfrak{p}=\mathfrak{q}\cap A$
($0_{\mathrm I}$, 6.3.1); basta, por tanto, en virtud de (17.1.10), probar que
$B_{\mathfrak{q}}/\mathfrak{p}B_{\mathfrak{q}}$ es regular, y como este anillo
es un anillo local en un ideal primo de
$A_{\mathfrak{p}}[T]/\mathfrak{p}A_{\mathfrak{p}}[T]=k[T]$ (donde $k$ es el
cuerpo residual $A_{\mathfrak{p}}/\mathfrak{p}A_{\mathfrak{p}}$), basta
probar que $C=k[T]$ es regular; ahora bien, como $C$ es un dominio de ideales
principales, los anillos locales en los ideales primos de $C$ son anillos de
valoración discreta o un cuerpo (para el ideal $(0)$), luego son regulares
(17.1.4), lo que termina la demostración.
\end{proof}

\begin{corollary}[17.3.8]
\label{0.17.3.8-es}
Si $A$ es un anillo regular, todo anillo de series formales
$A[[T_1,\ldots,T_n]]$ es regular.
\end{corollary}

\begin{proof}
Sea $\mathfrak{J}$ el ideal engendrado por los $T_i$ en el anillo de polinomios
$A[T_1,\ldots,T_n]$; como este último es regular por (17.3.7), y
$A[[T_1,\ldots,T_n]]$ es el completado de $A[T_1,\ldots,T_n]$ para la
topología $\mathfrak{J}$-preádica (Bourbaki, \emph{Alg. comm.}, cap.~III,
\S~2, n\textsuperscript{o}~12, \emph{Ejemplo 1}), la conclusión resulta del
\end{proof}

\begin{lemma}[17.3.8.1]
\label{0.17.3.8.1-es}
Sean $A$ un anillo regular, $\mathfrak{J}$ un ideal de $A$,
$\widehat{A}$ el completado separado de $A$ para la topología
$\mathfrak{J}$-preádica. Entonces $\widehat{A}$ es regular.
\end{lemma}

\begin{proof}
Basta, en efecto, (17.3.6) ver que, para todo ideal máximo
$\mathfrak{n}$ de $\widehat{A}$, el anillo local
$\widehat{A}_{\mathfrak{n}}$ es regular; se sabe (Bourbaki,
\emph{Alg. comm.}, cap.~III, \S~3, n\textsuperscript{o}~4, prop.~8) que
$\mathfrak{n}$
\oldpage[0\textsubscript{IV-1}]{51}
es de la forma $\mathfrak{m}\widehat{A}$, donde $\mathfrak{m}$ es un ideal
máximo de $A$ que contiene a $\mathfrak{J}$, y que el homomorfismo canónico
$A\to\widehat{A}$ da un homomorfismo \emph{inyectivo}
$A_{\mathfrak{m}}\to\widehat{A}_{\mathfrak{n}}$, tal que la topología
$\mathfrak{n}\widehat{A}_{\mathfrak{n}}$-preádica de
$\widehat{A}_{\mathfrak{n}}$ induzca sobre $A_{\mathfrak{m}}$ la topología
$\mathfrak{m}A_{\mathfrak{m}}$-preádica, y que $A_{\mathfrak{m}}$ sea denso
en $\widehat{A}_{\mathfrak{n}}$ para la topología
$\mathfrak{n}\widehat{A}_{\mathfrak{n}}$-preádica; por consiguiente, los
\emph{completados} de los anillos locales noetherianos $A_{\mathfrak{m}}$ y
$\widehat{A}_{\mathfrak{n}}$ son canónicamente \emph{isomorfos}. Ahora bien,
por hipótesis $A_{\mathfrak{m}}$ es regular, luego también lo es su completado
(17.1.5), y como el completado de $\widehat{A}_{\mathfrak{n}}$ es regular,
también lo es $\widehat{A}_{\mathfrak{n}}$ por (17.1.5).
\end{proof}

\begin{corollary}[17.3.9]
\label{0.17.3.9-es}
Sea $A$ un anillo noetheriano, cociente de un anillo noetheriano regular $B$.
Si $C$ es una $A$-álgebra de tipo finito, todo anillo de fracciones $S^{-1}C$
de $C$ es cociente de un anillo regular.
\end{corollary}

\begin{proof}
En efecto, $C$ es cociente de un anillo de polinomios $A[T_1,\ldots,T_n]$;
como $A=B/\mathfrak{J}$, donde $\mathfrak{J}$ es un ideal de $B$, se tiene
$A[T_1,\ldots,T_n]=B[T_1,\ldots,T_n]/\mathfrak{J}B[T_1,\ldots,T_n]$; en
virtud de (17.3.7), se puede, por tanto, limitar al caso en que $C$ es cociente
de un anillo regular $B'$; pero si $S'$ es la imagen inversa de $S$ en $B'$,
$S'$ es un conjunto multiplicativamente cerrado de $B'$ y $C$ es un anillo
cociente de $S'^{-1}B'$, de modo que finalmente todo se reduce a ver que,
cuando $C$ es \emph{regular}, también lo es $S^{-1}C$, lo que se ha visto en
(17.3.6).
\end{proof}

La importancia de los cocientes de anillos regulares se debe, entre otras
cosas, a la propiedad precedente y al hecho de que son anillos
\emph{catenarios} (16.5.12).

Todos los anillos importantes en las aplicaciones a la Geometría algebraica
son cocientes de anillos regulares.
